Online harassment on social media platforms poses a growing threat to user safety, and the volume of user-generated content makes manual moderation impractical at scale. Existing detectors often rely on a single text representation---lexical n-grams, sequential structure, or contextual embeddings---and therefore miss complementary cues present in abusive posts. This paper presents TriFuse, an end-to-end multi-view classifier that jointly models lexical, semantic, and structural text features within one trainable network. A CNN branch captures phrase-level patterns, a Transformer encoder branch models global context, and a BiLSTM branch encodes sequential discourse; cross-view interaction and temperature-scaled attention fusion combine these views at the representation level rather than through late ensemble voting. TriFuse is evaluated on the Shahane multi-platform corpus (Dataset~I, 229{,}228 instances) and the Davidson hate-speech corpus (Dataset~II, 24{,}783 instances) under matched preprocessing, splits, and training budgets. On Dataset~I, TriFuse obtains 94.32\% accuracy and 94.12\% F1-score on the held-out test set, and 94.29\,$\pm$\,0.13\% accuracy under stratified 5-fold cross-validation, outperforming Random Forest, LightGBM, CNN, BiLSTM, and Tuned LSTM. On Dataset~II, it reaches 95.87\% accuracy, 95.91\% F1-score, and 95.37\,$\pm$\,0.23\% cross-validated accuracy with the same ranking. An ablation study on both datasets shows that the full architecture outperforms single-branch, pairwise, and late-fusion configurations. Reported gains over the strongest neural baselines are incremental but consistent across folds and corpora; limitations regarding statistical testing, scope, and computational cost are discussed in Section~VII.
